% Paper 6 (T6) -- IEEE JBHI special issue "Integrating Imaging with Personalized Cardiovascular Medicine"
% Deadline 2026-10-15. Target: 8 pages (over-length charges start at page 9).
% Plan and word budgets: BENCHMARK-AND-PLAN-JBHI-2026-10-03.md
\documentclass[journal]{IEEEtran}
\usepackage{amsmath,amssymb}
\usepackage{graphicx}
\usepackage{booktabs}
\usepackage{siunitx}
\usepackage[hidelinks]{hyperref}

\begin{document}

\title{Boundary-Condition Tuning Conceals Topological Segmentation Error in Reduced-Order Coronary FFR}
% Working title; alternative from STUDY-PLAN-v2 sec. 0:
% "Matched at the Outlet, Wrong in the Tree: ..."

\author{Fadillah~Yamin,
        Mohd~Azan~Mohammed~Sapardi,
        Ming~Kwang~Tan,
        Xin~Wang,
        and~Mohd-Zulhilmi~Paiz~Ismadi%
\thanks{F. Yamin, X. Wang and M.-Z. P. Ismadi are with the Department of Mechanical Engineering, School of
Engineering, Monash University Malaysia, 47500 Bandar Sunway, Selangor, Malaysia (corresponding author:
M.-Z. P. Ismadi, e-mail: mohd.zulhilmi@monash.edu).}%
\thanks{M. A. Mohammed Sapardi is with the Department of Mechanical Engineering, Kulliyyah of Engineering,
International Islamic University Malaysia, 53100 Kuala Lumpur, Malaysia.}%
\thanks{M. K. Tan is with the Department of Mechanical Engineering, School of Engineering, Monash University
Malaysia, and the Department of Mechanical Engineering, Ming Chi University of Technology, New Taipei City 243,
Taiwan.}}

\markboth{IEEE Journal of Biomedical and Health Informatics}{Yamin \MakeLowercase{\textit{et al.}}: BC Tuning Conceals Segmentation Error in Coronary FFR}

\maketitle

\begin{abstract}
% <= 250 words, one paragraph, no abbreviations, no references or equations (JBHI rule).
% Context (1) -> gap (1) -> aim (1) -> methods (2-3) -> results with numbers (2-4) -> conclusion (1) -> significance (1).
[ABSTRACT -- written last, after the results.]
\end{abstract}

\begin{IEEEkeywords}
Boundary conditions, coronary artery disease, fractional flow reserve, image segmentation, reduced-order modeling,
uncertainty.
\end{IEEEkeywords}

% ---------------------------------------------------------------------------------------------------------------
\section{Introduction}
\IEEEPARstart{F}{ractional} flow reserve (FFR) at or below 0.80 identifies a coronary stenosis for revascularisation
\cite{tonino2009}. Computed tomography (CT)-derived FFR estimates the same quantity without a pressure wire, by
solving a flow model on a lumen segmented from coronary CT angiography \cite{taylor2013,norgaard2014,itu2016}. The
model has two image-dependent inputs: the lumen geometry, which comes from segmentation, and the outlet boundary
conditions, which are derived from vessel calibre through allometric scaling and represent the microvascular bed
\cite{kim2010,murray1926}. An error in either input propagates to the computed FFR.

Published studies quantify the sensitivity of CT-derived FFR to lumen-radius uncertainty
\cite{sankaran2015tmi,sankaran2015cmame,fernandez2024} and compare the influence of geometry with that of boundary
conditions in coronary, pulmonary and cerebral vessels \cite{sankaran2016,tanade2022,colebank2019,korte2023}. These
studies perturb lumen calibre with the branching structure held fixed, and they report continuous sensitivity rather
than changes in the treatment decision. Automatic segmentation, however, also produces topological errors. In a
recent coronary benchmark, vessel breaks were present in the predictions of every evaluated method despite high
overlap scores \cite{bransby2026}. A missed side branch or a broken vessel removes part of the microvascular bed from
the model, which changes both the flow through the stenosis and the boundary conditions derived from the tree.

Model pipelines adjust boundary conditions so that predicted flow matches a physiological target, and this
adjustment can compensate for a geometric discrepancy. Re-tuning the microvascular resistance absorbed a change in
modelled side-branch flow with little change in diagnostic accuracy \cite{gosling2020}. Geometry and boundary
conditions have been adjusted jointly to match the same perfusion targets \cite{menon2024}. A change of the outlet
boundary-condition model left the area under the receiver operating characteristic curve unchanged (0.845) while
sensitivity changed from 58.1\% to 68.6\% \cite{fossan2026}. With outlet resistance held fixed, removal of a side
branch downstream of a stenosis changed FFR by 13--15\%, an effect that depends on whether resistance is re-tuned
\cite{gamage2022}. A model that matches its perfusion targets after re-tuning can therefore still carry a
segmentation error into the computed FFR. Whether this changes the decision at 0.80, and how often, has not been
tested under controlled conditions.

This study tests it. We inserted controlled stenoses into 150 real coronary trees from a public CT angiography
dataset, applied four segmentation error types, and recomputed FFR under three boundary-condition protocols: fixed,
re-derived from the corrupted tree, and re-tuned to the clean tree's territory perfusion. The design and analysis
were pre-registered. The contributions are: (1) a controlled ablation of boundary-condition protocols against
topological and calibre segmentation errors, with the 0.80 decision as the outcome and a physiological-noise floor
as the comparator; (2) the proportion of models that pass a perfusion-based validation check while their FFR is
materially wrong; and (3) a three-dimensional computational fluid dynamics (CFD) case study of the same mechanism on
a segmented lumen.

% ---------------------------------------------------------------------------------------------------------------
\section{Methods}
\subsection{Data and Cohort}
ImageCAS-X provides voxel-wise lumen and coronary-segment annotations, centerlines and surfaces for 800 CT angiography
scans of the ImageCAS dataset \cite{zeng2023,bransby2026,imagecasx_data}. We used its 160-scan test split, in which
each scan was re-annotated by a second analyst. Inter-observer agreement on this split was a Dice similarity
coefficient of 92.8\% and a 95th-percentile Hausdorff distance (HD95) of 2.46~mm \cite{bransby2026}. Each scan gives a
left and a right coronary tree. The local radius at each centerline point is the Euclidean distance transform of the
lumen mask, less half a voxel.

The natural cohort contains few lesions near the decision threshold, so stenoses were inserted into healthy vessels.
A host slot is a main vessel (left anterior descending, left circumflex or right coronary artery), a location
(20 or 45~mm from the vessel origin) and a length $L$ (10 or 20~mm). Eligible hosts had a healthy reference radius
$r_\mathrm{fit} \geq 1.0$~mm at the lesion centre, native narrowing below 40\% diameter stenosis (DS), a resolved
measurement point 20~mm distal to the lesion and a pre-insertion FFR of at least 0.90. The lesion is an axisymmetric
cosine narrowing,
\begin{equation}
r'(s) = \min\!\left(r,\; (1-w)\,r + w\,r_\mathrm{fit}\,(1-\mathrm{DS})\right),\;
w = \tfrac{1}{2}\left[1+\cos\tfrac{2\pi (s-c)}{L}\right]
\end{equation}
for $|s-c|<L/2$, with severities of 40--80\% DS. The 6\,944 eligible instances were stratified into six bands of
baseline FFR (0.65--0.95), and 150 instances were selected (50 per vessel) from 108 trees of 93 patients. The
selected instances form a frozen, hashed list.

\subsection{Reduced-Order FFR Model}
The tree is a network with one node per centerline point. Each element has a Poiseuille resistance
$8\mu\,\Delta s/(\pi r^4)$ on the actual radius, and each stenosis throat adds an expansion loss
$K Q|Q|$ with $K = \rho K_t (A_0/A_s - 1)^2 / (2A_0^2)$ and $K_t = 1.52$, where $A_0$ and $A_s$ are the reference and
throat areas. The microvascular bed connects nodes to venous pressure through conductances proportional to a healthy
reference radius $r_\mathrm{ref}$, scaled by one bed constant $C$. Two bed structures were analysed. In the leaky
bed, flow leaves along the vessel wall according to Murray's law, with conductance $(r_\mathrm{ref}^3 - \sum
r_{\mathrm{ref},\mathrm{child}}^3)/C$ at each node \cite{murray1926}. In the discrete bed, flow leaves only at outlets
of the tree truncated at $r_\mathrm{ref} = 0.60$~mm, with conductance $r_\mathrm{ref}^{2.66}/C$, where 2.66 is the
myocardial mass--diameter exponent \cite{choy2008}. The discrete bed is the structure a three-dimensional model can
represent. $C$ was calibrated on the healthy-equivalent tree so that inflow equals the hyperaemic demand $Q = k
r_\mathrm{in}^3$ with $k = 562$~s$^{-1}$. Aortic and venous pressures were 90 and 5~mmHg, blood viscosity
$\mu = 0.004$~Pa\,s and density $\rho = 1060$~kg\,m$^{-3}$. FFR is the pressure at the measurement point divided by
aortic pressure. The bed constant and reference radii are computed from the original tree, so an inserted lesion
changes only the epicardial radius.

\subsection{Segmentation Error Types}
\begin{table}[!t]
\caption{Segmentation Error Types and Boundary-Condition Protocols}
\label{tab:design}
\centering\footnotesize
\begin{tabular}{@{}p{0.9cm}p{3.6cm}p{2.9cm}@{}}
\toprule
 & Operation & Magnitude (source)\\
\midrule
T1 & delete largest side branch distal to lesion & largest eligible (design)\\
T2 & truncate host vessel beyond lesion & 25~mm distal (design)\\
T3 & lengthen lesion & +2.46~mm (HD95)\\
T4 & scale radius from lesion distally & $\times$0.930 (Dice 92.8\%)\\
\midrule
A & clean bed conductances and $C$ & fixed\\
B & bed rule and $C$ on corrupted tree & re-derived\\
C & one scaling of $C$ to territory flows & fitted\\
\bottomrule
\end{tabular}
\end{table}
Four error types were applied to each instance (Table~\ref{tab:design}). T1 (missed side branch) deletes the largest
side branch leaving the host vessel distal to the lesion, with its subtree. T2 (vessel break) truncates the host
vessel 25~mm beyond the distal edge of the lesion, so the measurement point remains interior. T3 (lesion length)
lengthens the lesion by 2.46~mm, the inter-observer HD95. T4 (taper) scales the radius by 0.930 from the proximal
edge of the lesion through all descendants, the coaxial-cylinder radius ratio equivalent to a Dice coefficient of
92.8\%. T3 and T4 derive from the measured inter-observer statistics, which the dataset authors describe as an upper
bound on agreement. T1 and T2 have no measured magnitude and are design choices. For T3 and T4 the set of modelled
nodes is fixed to that of the clean tree.

\subsection{Boundary-Condition Protocols}
Each corrupted tree was solved under three protocols. In Protocol A (fixed), surviving nodes keep the clean tree's
bed conductances and $C$, and the bed of deleted vessels is lost. In Protocol B (re-derived), the bed rule is applied
to the corrupted tree and $C$ is recalibrated to its own hyperaemic demand, as a deployed pipeline would do. In
Protocol C (flow-matched), one global scaling of $C$ is fitted so that the flow of each perfusion territory matches the
clean tree's full territory flow, including the share of any deleted branch. Territories are the subtrees at the
first bifurcation of the clean tree. The fit has one free parameter and at least two targets, and it uses a 61-point
scan over $\pm 1.5$ decades followed by bounded refinement. A fit at a search bound is recorded as a failure. Protocol
C matches flow only: matching both flow and pressure at the outlets of a steady tree reproduces the clean FFR by
construction.

\subsection{Three-Dimensional Case Study}
One instance, a 20~mm lesion of 80\% DS in the proximal left anterior descending artery, was solved in three
dimensions in its clean, baseline and T1 forms. Lumens were built from the edited masks by marching cubes and Taubin
smoothing, with the lesion applied by vertex deformation and the side branch removed by mask editing. Meshes of
3.5--3.9 million cells (cfMesh, 25~\textmu m cells within $\pm 4$~mm of the throat, four boundary layers) were solved
for steady, laminar, Newtonian flow with OpenFOAM (ESI v2406). The discretisation uncertainty in FFR, from a grid
convergence study on an idealised 70\% stenosis, was $5.5\times10^{-4}$. Outlets received either the clean tree's
resistances (Protocol A) or prescribed flows equal to the clean full-territory flows (a per-outlet form of
Protocol C). A reduced-order twin of each geometry, with the radius replaced by the area-equivalent radius of the
meshed lumen and identical outlet conditions, isolates the effect of model fidelity from that of radius definition.

\subsection{Outcomes and Statistical Analysis}
For each instance, error type, protocol and bed, the change $\Delta$FFR is measured against the instance's own clean
baseline in the same bed, and a flip is a change of classification at 0.80. The perfusion residual is the
root-mean-square relative error of the territory flows. A model passes validation when the residual is below 10\%,
which is stricter than the 95\% bound of 13--16\% implied by the repeatability of regional perfusion measurement
\cite{schindler2023,brown2018}, and
is materially wrong when $|\Delta\mathrm{FFR}| > 0.05$.

Flips were modelled by mixed-effects logistic regression with error type, protocol and their interaction, severity,
length and baseline band as fixed effects and a patient-level random effect with a nested slot term. Protocol
contrasts are paired within instance and tested with McNemar's exact test. Proportions are reported with Wilson 95\%
intervals. The Holm--Bonferroni procedure controls the family-wise error rate at 0.05 across the six registered
hypotheses. A claim is made only where direction agrees in both bed structures. Each flip rate is compared with a
physiological-noise floor from repeat invasive FFR (standard deviation 0.018) \cite{johnson2015,petraco2013} and with
a Monte Carlo simulation of within-subject physiological variation on the clean anatomy.

The protocol, analysis plan, cohort and code were registered before the ablation was run (OSF:
\textit{[DOI]}; code: \textit{[DOI]}). The registration includes hypotheses on 0D--3D concordance and on a
deployment-time detector; their results are not part of this report. Deviations from the registered plan are listed
in the Supplementary Material.

% ---------------------------------------------------------------------------------------------------------------
\section{Results}
% ~1,300 words. Subsections mirror the evaluation questions.
\subsection{Decision Changes by Error Type and Protocol}      % H1, H3; Fig. 2
\subsection{Absorption Under Tuned Boundary Conditions}       % H2; Fig. 3, Table II
\subsection{Side-Branch Loss Under Fixed and Tuned Resistance} % H6; Fig. 4
\subsection{Three-Dimensional Case Study}
% Data in hand (returns/2026-09-26; results/cfd_M1/M1_scan14_0D_vs_3D-2026-10-03.csv). M1 D7 re-mesh test pending.
With the clean tree's resistances at the outlets, removal of the side branch raised the 3D FFR at the measurement
point from 0.870 to 0.944 ($\Delta$FFR $= +0.074$). With the clean territory flows prescribed at the outlets, the same
error changed FFR by $-0.0007$ (0.892 in both geometries). The
reduced-order twins gave $+0.081$ and $-0.0005$ for the same two conditions, and their baseline FFR was within 0.011
of the 3D values (Fig.~\ref{fig:case3d}).
\begin{figure}[!t]\centering
\fbox{\parbox[c][4cm][c]{0.9\columnwidth}{\centering [Fig.~5: pressure along the LAD, baseline and T1, fixed
resistances and prescribed flows; 3D markers, reduced-order twin lines.]}}
\caption{Three-dimensional case study (proximal left anterior descending artery, 80\% DS).}
\label{fig:case3d}
\end{figure} The meshed lumen was wider than the radius used by the reduced-order model
(median difference 0.14~mm; throat area-equivalent radius 0.276 against 0.231~mm), and on the requested radius the
reduced-order baseline was 0.761.


% ---------------------------------------------------------------------------------------------------------------
\section{Discussion}
% ~900 words.
\subsection{Concealment at Calibration}
\subsection{Relation to Published Studies}
\subsection{Implications for Model Validation}
\subsection{Limitations}                  % ~200 words

% ---------------------------------------------------------------------------------------------------------------
\section{Conclusion}
% ~100 words. Findings only.

\section*{Data and Code Availability}
ImageCAS-X is available under CC BY 4.0 (doi:10.5281/zenodo.21887809). Code and the frozen cohort: \textit{[DOI]}.
Pre-registration: \textit{[OSF DOI]}.

\section*{Acknowledgment}

\bibliographystyle{IEEEtran}
\bibliography{refs}

\end{document}
